
\documentclass{article}
\usepackage{colortbl}
\usepackage{makecell}
\usepackage{multirow}
\usepackage{supertabular}

\begin{document}

\newcounter{utterance}

\centering \large Interaction Transcript for game `wordle\_withclue', experiment `medium\_frequency\_words\_clue\_no\_critic', episode 10 with reflexion{-}Qwen3.5{-}27B.
\vspace{24pt}

{ \footnotesize  \setcounter{utterance}{1}
\setlength{\tabcolsep}{0pt}
\begin{supertabular}{c@{$\;$}|p{.15\linewidth}@{}p{.15\linewidth}p{.15\linewidth}p{.15\linewidth}p{.15\linewidth}p{.15\linewidth}}
    \# & \multicolumn{2}{c}{Player} && \multicolumn{2}{c}{Game Master} \\
    \hline

    \theutterance \stepcounter{utterance}  
    & & & \multicolumn{4}{p{0.6\linewidth}}{
        \cellcolor[rgb]{0.9,0.9,0.9}{
            \makecell[{{p{\linewidth}}}]{
                \texttt{\tiny{[P1$\langle$GM]}}
                \texttt{You are a language wizard who likes to guess words by using the given rules.} \\
\\ 
\texttt{Welcome to Wordle! You have six attempts to guess the target word, a valid English word of five lowercase letters (a{-}z). Please use the tags "explanation:" and "guess:" to provide a concise explanation for each guess.} \\
\\ 
\texttt{To help you make an informed guess, you will receive a clue for the word, such as} \\
\texttt{clue: snowy white.} \\
\\ 
\texttt{Here is an example guess based on the clue:} \\
\texttt{explanation: In the fairy tail Snow White, the girl is killed because she eats a poisoned apple. And the word apple has 5 letters.} \\
\texttt{guess: apple} \\
\\ 
\texttt{After each guess, your answer will be validated, and you will receive feedback indicating which letters are correct (green), which letters are correct but in the wrong position (yellow), and which letters are incorrect (red). This feedback can be useful in determining which letters to include or exclude in your next guess.} \\
\\ 
\texttt{For example, the feedback for "apple" might be:} \\
\texttt{guess\_feedback: a<yellow> p<yellow> p<green> l<yellow> e<red>} \\
\\ 
\texttt{The explanation should contain details about how both the guess\_feedback and clue work together in tandem to formulate a new guess.} \\
\\ 
\texttt{Do not generate any other text.} \\
\\ 
\texttt{Let's begin with your first guess.} \\
\\ 
\texttt{clue: tithing} \\
            }
        }
    }
     \\ \\

    \theutterance \stepcounter{utterance}  
    & \multicolumn{4}{p{0.6\linewidth}}{
        \cellcolor[rgb]{0.561,0.737,0.561}{
            \makecell[{{p{\linewidth}}}]{
                \texttt{\tiny{[P1$\rangle$GM]}}
                \texttt{explanation: The clue "tithing" refers to the practice of giving a tenth of one's income, often to a church. A common 5{-}letter word associated with this concept is "tithe," which means the act of giving a tenth.} \\
\texttt{guess: tithe} \\
            }
        }
    }
    & & \\ \\

    \theutterance \stepcounter{utterance}  
    & & & \multicolumn{2}{p{0.3\linewidth}}{
        \cellcolor[rgb]{0.9,0.9,0.9}{
            \makecell[{{p{\linewidth}}}]{
                \texttt{\tiny{[GM$|$GM]}}
                \texttt{attempts = 1} \\
\texttt{target\_word = tenth} \\
\texttt{target\_word\_clue = tithing} \\
\texttt{guess = tithe} \\
\texttt{guess\_feedback = t<green> i<red> t<yellow> h<yellow> e<yellow>} \\
            }
        }
    }
    & & \\ \\

    \theutterance \stepcounter{utterance}  
    & & & \multicolumn{4}{p{0.6\linewidth}}{
        \cellcolor[rgb]{0.9,0.9,0.9}{
            \makecell[{{p{\linewidth}}}]{
                \texttt{\tiny{[P1$\langle$GM]}}
                \texttt{guess\_feedback: t<green> i<red> t<yellow> h<yellow> e<yellow>} \\
\\ 
\texttt{Provide your response only in this format.} \\
\texttt{explanation: details} \\
\texttt{guess: word} \\
            }
        }
    }
     \\ \\

    \theutterance \stepcounter{utterance}  
    & \multicolumn{4}{p{0.6\linewidth}}{
        \cellcolor[rgb]{0.561,0.737,0.561}{
            \makecell[{{p{\linewidth}}}]{
                \texttt{\tiny{[P1$\rangle$GM]}}
                \texttt{explanation: The feedback shows 't' is correct in the first position (green). 'i', 'h', and 'e' are in the word but in wrong positions (yellow), while the second 't' is also misplaced. Since 't' is already fixed at the start, the second 't' must be in position 3, 4, or 5. The word "tithing" relates to giving a tenth, and "tithe" was the guess. With 't' at pos 1, and 'i', 'h', 'e' needing to move, and another 't' needed, a strong candidate is "thief" (no, 'f' not in feedback), "thine" (no 'n'), "thyme" (no 'm'). Wait, the letters available are T, I, H, E and another T. The word must be 5 letters. The clue is "tithing". The word "tithe" has T, I, T, H, E. Feedback: T(green), I(yellow), T(yellow), H(yellow), E(yellow). This means the word starts with T. The other letters I, T, H, E are in the word but not in positions 2, 3, 4, 5 respectively. So position 2 is not I, position 3 is not T, position 4 is not H, position 5 is not E. We need to place I, T, H, E in positions 2, 3, 4, 5 such that none are in their original spots from "tithe". Original: 1:T, 2:I, 3:T, 4:H, 5:E. New: 1:T. Remaining slots 2,3,4,5 for \{I, T, H, E\}. Constraints: Pos2 != I, Pos3 != T, Pos4 != H, Pos5 != E. Possible permutations: T{-}H{-}I{-}E{-}T? No, T is used. Letters are T, I, T, H, E. So we have two T's. One is at pos 1. The other T must be in 2, 3, 4, or 5. But Pos3 != T. So T is in 2, 4, or 5. Let's try "thite"? No. "theit"? No. "theti"? No. Maybe the word is "thine"? No, 'n' not in feedback. "thyme"? No. "thief"? No. "thine" is 5 letters but 'n' is not in the feedback. The feedback only lists letters in the guess. If the target word had 'n', it would be red or yellow. Since 'n' wasn't guessed, we don't know. But} \\
            }
        }
    }
    & & \\ \\

    \theutterance \stepcounter{utterance}  
    & & & \multicolumn{2}{p{0.3\linewidth}}{
        \cellcolor[rgb]{0.9,0.9,0.9}{
            \makecell[{{p{\linewidth}}}]{
                \texttt{\tiny{[GM$|$GM]}}
                \texttt{Error: The response should contain the 'guess:' keyword exactly once.} \\
            }
        }
    }
    & & \\ \\

    \theutterance \stepcounter{utterance}  
    & & & \multicolumn{2}{p{0.3\linewidth}}{
        \cellcolor[rgb]{0.9,0.9,0.9}{
            \makecell[{{p{\linewidth}}}]{
                \texttt{\tiny{[GM$|$GM]}}
                \texttt{game\_result = ABORT} \\
            }
        }
    }
    & & \\ \\

    \theutterance \stepcounter{utterance}  
    & & \multicolumn{4}{p{0.6\linewidth}}{
        \cellcolor[rgb]{0.95,0.95,0.95}{
            \makecell[{{p{\linewidth}}}]{
                \texttt{\tiny{[P1$\leftrightarrow$P]}}
                \texttt{Post{-}episode Reflexion (not sent to GM):} \\
\texttt{New plan: The previous attempt failed because I violated the instruction "Do not generate any other text" by including a long, rambling internal monologue and logical deduction within the "explanation" field, rather than providing a concise summary. Additionally, I failed to provide a valid "guess:" line with a 5{-}letter word, leaving the response incomplete. Next time, I will strictly adhere to the format by providing only a brief explanation of how the feedback and clue inform the next guess, followed immediately by a valid 5{-}letter word guess. I will perform all logical deductions internally before generating the response to ensure the output contains no extraneous text or incomplete sentences.} \\
            }
        }
    }
    & \\ \\

\end{supertabular}
}

\end{document}
